\documentclass[10pt,a4paper]{amsart}
\usepackage[margin=19mm]{geometry}
\usepackage{amsmath,amssymb,mathtools}
\usepackage[scr=rsfs,cal=euler]{mathalfa}
\usepackage{xfrac}
\usepackage[unicode,hidelinks,bookmarks=false]{hyperref}
\newcommand{\ie}{i.e.\ }
\newcommand{\cf}{cf.\ }
\newcommand{\resp}{respectively\ }
\newcommand{\Subsection}{\subsection}
\providecommand{\nobreakdash}{\nobreak\hbox{-}}
\setlength{\emergencystretch}{2em}
\multlinegap0pt
\usepackage{amssymb}
\usepackage{tikz}
\newtheorem{lem}{Lemma}
\title{On the Fitting ideals of anticyclotomic Selmer groups of elliptic curves with good ordinary reduction}
\author{Chan-Ho Kim}
\date{Source: arXiv:2505.09125v2 (CC BY 4.0)}
\begin{document}
\maketitle

\noindent\emph{Source and licence.} On the Fitting ideals of anticyclotomic Selmer groups of elliptic curves with good ordinary reduction. Version arXiv:2505.09125v2 (CC BY 4.0), distributed by arXiv under the Creative Commons Attribution 4.0 International licence, http://creativecommons.org/licenses/by/4.0/, as recorded in the arXiv licence metadata for this exact version and shown on the abstract page https://arxiv.org/abs/2505.09125v2. Full licence notice: \texttt{licenses/arxiv-2505.09125-CC-BY-4.0.txt}.\par\smallskip

\noindent\textit{Editorial scope.} Counted unit: the article's lem lem:reduction-generators (source0.tex lines 312--315). The article's complete original proof of it is delivered with it (source0.tex lines 316--333, one \texttt{proof} environment of 18 lines). No mathematics is written, repaired or re-derived: the statement and the proof are the article's own lines, in the article's own notation, and no retained line is substituted or re-flowed.  Retained uncounted as the source-local context the proof needs: lines 294--296.  Nothing was omitted from any retained span, so the package carries no omission record.  No retained line contains a citation, so the wrapper carries no bibliography and no bibliography entry.  No retained line contains a figure environment, an image-inclusion command or drawing code, so no image is delivered and the package declares no asset dependency.  Editorial formatting only, in addition: an AMS \texttt{amsart} preamble that restates the article's own declarations and macros used by the retained lines, namely the environments \texttt{lem} on separate counters, together with the standard packages those lines need.  The retained environments print in this document as Lemma 1 in order of appearance, whereas the article prints the same items with the numbering of its own full document.  The complete native source of the article as distributed in the arXiv e-print for this exact version is bundled unchanged as \texttt{sources/178095/source0.tex}.
\begin{equation} \label{eqn:three-term-relation}
\pi_{n+1, n} \left( \theta(E/K_{n+1}) \right) = a_p \cdot \theta(E/K_n) - \nu_{n-1, n} \left( \theta(E/K_{n-1}) \right)
\end{equation}

\begin{lem} \label{lem:reduction-generators}
We have an equality of ideals of $\Lambda_n$
$$\left( \theta(E/K_n) , \nu_{n-1, n} \left( \theta(E/K_{n-1}) \right) \right) = \left( \nu_{m, n} \left( \theta(E/K_{m}) \right) : 0 \leq m \leq n \right) .$$
\end{lem}

\begin{proof}
From the three term relation (\ref{eqn:three-term-relation}), we have
$$\nu_{n-1, n} \left( \pi_{n, n-1} \left( \theta(E/K_{n}) \right) \right) = a_p \cdot \nu_{n-1, n}  \left( \theta(E/K_{n-1}) \right) - \nu_{n-2, n} \left( \theta(E/K_{n-2}) \right) $$
for $n \geq 2$.
Since 
$\nu_{n-1, n} \left( \pi_{n, n-1} \left( \theta(E/K_{n}) \right) \right) = f_n \cdot \theta(E/K_{n})$
for some $f_n \in \Lambda_n$, we have
$$\nu_{n-2, n} \left( \theta(E/K_{n-2}) \right) \in \left( \theta(E/K_n) , \nu_{n-1, n} \left( \theta(E/K_{n-1}) \right) \right) \subseteq \Lambda_n.$$
In the same manner, we can obtain
$$\nu_{n-3, n-1} \left( \theta(E/K_{n-3}) \right) \in \left( \theta(E/K_{n-1}) , \nu_{n-2, n-1} \left( \theta(E/K_{n-2}) \right) \right) \subseteq \Lambda_{n-1}. $$
By taking $\nu_{n-1, n}$, we have
\begin{align*}
\nu_{n-3, n} \left( \theta(E/K_{n-3}) \right) & \in \left( \nu_{n-1, n} \left( \theta(E/K_{n-1}) \right) , \nu_{n-2, n} \left( \theta(E/K_{n-2}) \right) \right) \\
& \subseteq \left( \theta(E/K_{n}), \nu_{n-1, n} \left( \theta(E/K_{n-1}) \right)  \right) \\
& \subseteq \Lambda_{n}. 
\end{align*}
By applying this argument recursively, the conclusion follows.
\end{proof}

\end{document}
